\documentclass[aspectratio=169]{beamer}

\usetheme{Madrid}
\usecolortheme{default}

\usepackage[utf8]{inputenc}
\usepackage[T1]{fontenc}
\usepackage[brazilian]{babel}
\usepackage{lmodern}
\usepackage{tikz}
\usepackage{booktabs}
\usepackage{listings}
\usepackage{xcolor}

\usetikzlibrary{positioning}

\definecolor{codegray}{RGB}{245,245,245}
\definecolor{codeblue}{RGB}{35,88,140}

\lstdefinestyle{proto}{
  basicstyle=\ttfamily\scriptsize,
  backgroundcolor=\color{codegray},
  keywordstyle=\color{codeblue}\bfseries,
  showstringspaces=false,
  columns=fullflexible,
  keepspaces=true,
  frame=single,
  breaklines=true
}

\title[Sistema Distribuido de Cidade Inteligente]{Sistema Distribuido de Cidade Inteligente}
\subtitle{Detalhes de Implementacao, Mensagens e Tecnologias}
\author{Distribuicao de Dados}
\institute{Universidade Federal do Ceara}
\date{2026.1}

\begin{document}

\begin{frame}
  \titlepage
\end{frame}

\begin{frame}{Objetivo do Sistema}
  \begin{itemize}
    \item Implementar um sistema distribuido para coleta, controle e consulta de dados urbanos.
    \item Integrar tres fontes de dados: clima, qualidade do ar e poste inteligente.
    \item Centralizar descoberta, leituras, historico, controle e consultas em um Gateway.
    \item Disponibilizar uma interface web analitica para visualizacao e operacao do sistema.
    \item Usar mensagens padronizadas com Protocol Buffers entre os processos.
  \end{itemize}
\end{frame}

\begin{frame}{Partes do Trabalho}
  \begin{enumerate}
    \item Descoberta dinamica das fontes de dados pelo Gateway.
    \item Recebimento e armazenamento de leituras das fontes.
    \item Cliente web analitico para visualizacao, controle e consultas agregadas.
  \end{enumerate}

  \vspace{0.4cm}
  \begin{block}{Foco da implementacao}
    Formato das mensagens, comunicacao entre processos, linguagens, bibliotecas e protocolos usados.
  \end{block}
\end{frame}

\begin{frame}{Arquitetura Geral}
  \centering
  \resizebox{0.9\textwidth}{!}{%
  \begin{tikzpicture}[
    every node/.style={font=\small, align=center},
    box/.style={draw, rounded corners, minimum width=2.8cm, minimum height=0.95cm}
  ]
    \node[box, fill=green!10] (weather) at (-4.8, 2.0) {weather-sensor\\Python};
    \node[box, fill=green!10] (air) at (-4.8, 0) {air-quality-sensor\\Python};
    \node[box, fill=green!10] (lamp) at (-4.8, -2.0) {lamppost\\Python};

    \node[box, fill=blue!10] (gateway) at (0, 0) {Gateway\\Go};
    \node[box, fill=gray!12] (csv) at (0, -2.0) {Historico CSV\\readings.csv};
    \node[box, fill=orange!12] (flask) at (4.2, 0) {Cliente Web\\Flask};
    \node[box, fill=orange!5] (browser) at (7.8, 0) {Navegador\\HTML/JS};

    \draw[->, thick, blue] (gateway) to[bend left=18] (weather);
    \draw[->, thick, blue] (gateway) to[bend left=12] (air);
    \draw[->, thick, blue] (gateway) to[bend right=18] (lamp);

    \draw[->, dashed, thick] (weather) to[bend left=8] (gateway);
    \draw[->, dashed, thick] (air) to[bend left=8] (gateway);
    \draw[->, dashed, thick] (lamp) to[bend right=8] (gateway);

    \draw[<->, thick, red] (gateway.south west) .. controls (-1.2,-3.1) and (-3.8,-3.1) .. (lamp.south east);
    \draw[<->, thick] (gateway) -- (flask);
    \draw[<->, thick] (flask) -- (browser);
    \draw[<->, thick] (gateway) -- (csv);
  \end{tikzpicture}
  }

  \vspace{0.25cm}
  \begin{minipage}{0.9\textwidth}
    \scriptsize
    \textcolor{blue}{linha azul}: discovery UDP multicast \texttt{239.0.0.1:9999}\\
    linha tracejada: leituras UDP \texttt{:11000}\\
    \textcolor{red}{linha vermelha}: controle TCP do \texttt{lamppost} \texttt{:10002}\\
    linha preta: TCP + Protobuf \texttt{:12000}, HTTP \texttt{:5000} e persistencia CSV
  \end{minipage}
\end{frame}

\begin{frame}{Processos e Portas}
  \small
  \begin{table}
    \centering
    \begin{tabular}{lll}
      \toprule
      Processo & Linguagem & Papel \\
      \midrule
      Gateway & Go & Descoberta, estado, historico, leituras e TCP \\
      weather-sensor & Python & Sensor de Temperatura e umidade \\
      air-quality-sensor & Python & Sensor de CO2, particulados e indice de ar \\
      lamppost & Python & Fonte controlavel de luminosidade e energia\\
      Cliente Web & Python/Flask & Interface web analitica \\
      \bottomrule
    \end{tabular}
  \end{table}

  \vspace{-0.1cm}
  \begin{itemize}
    \item Discovery UDP multicast: \texttt{239.0.0.1:9999}
    \item Servidor UDP (Gateway): \texttt{11000}
    \item Servidor TCP (Gateway): \texttt{12000}
    \item Servidor de Controle TCP (Lamppost): \texttt{10002}
    \item Cliente Web Flask: \texttt{5000}
  \end{itemize}
\end{frame}

\begin{frame}{Parte 1: Descoberta de Fontes}
  \begin{itemize}
    \item O Gateway envia uma requisicao UDP multicast para o grupo de descoberta.
    \item Cada fonte responde por UDP usando \texttt{DiscoveryResponse} serializado com Protobuf.
    \item A resposta informa nome, tipo, endereco de controle e status operacional.
    \item O Gateway registra as fontes em um estado central protegido por \texttt{sync.RWMutex}.
    \item A descoberta periodica ocorre a cada 5 segundos.
    \item Fontes conhecidas que nao respondem a uma rodada sao marcadas como \texttt{OFFLINE}.
  \end{itemize}

\end{frame}

\begin{frame}[fragile]{Formato das Mensagens de Descoberta}
\begin{lstlisting}[style=proto]
message DiscoveryRequest {
  string gateway_id = 1;
  string gateway_ip = 2;
  int32 readings_port = 3;
  int32 client_port = 4;
}

message DiscoveryResponse {
  string source_name = 1;
  string source_type = 2;
  string ip = 3;
  int32 control_port = 4;
  bool controllable = 5;
  string status = 6;
}
\end{lstlisting}
\end{frame}

\begin{frame}{Parte 2: Envio de Leituras}
  \begin{itemize}
    \item As fontes enviam leituras periodicas ao Gateway por UDP.
    \item Cada pacote contem o nome da fonte, timestamp e uma leitura tipada.
    \item O campo \texttt{oneof} permite reaproveitar o mesmo envelope para diferentes sensores.
    \item O Gateway converte as leituras recebidas em metricas internas com valor, unidade e timestamp.
    \item As leituras sao mantidas em memoria e persistidas em arquivo CSV.
  \end{itemize}
\end{frame}

\begin{frame}[fragile]{Envelope de Leituras}
\begin{lstlisting}[style=proto]
message ReadingPacket {
  string source_name = 1;
  int64 timestamp_unix_ms = 2;

  oneof reading {
    temperature.TemperatureReading temperature = 3;
    airquality.AirQualityReading air_quality = 4;
    lamppost.LamppostReading lamppost = 5;
  }
}
\end{lstlisting}

  \begin{block}{Vantagem do \texttt{oneof}}
    O Gateway recebe um unico tipo de pacote, mas preserva o formato especifico de cada fonte.
  \end{block}
\end{frame}

\begin{frame}[fragile]{Formatos de Leituras por Fonte}
\begin{lstlisting}[style=proto]
message TemperatureReading {
  double temperature_celsius = 1;
  double humidity_percent = 2;
}

message AirQualityReading {
  double co2_ppm = 1;
  double particulate_matter_ug_m3 = 2;
  double air_quality_index = 3;
}

message LamppostReading {
  double luminosity_percent = 1;
  double energy_consumption_kwh = 2;
  bool light_on = 3;
}
\end{lstlisting}
\end{frame}

\begin{frame}{Metricas Armazenadas pelo Gateway}
  \begin{table}
    \centering
    \begin{tabular}{lll}
      \toprule
      Fonte & Metrica & Unidade \\
      \midrule
      weather & \texttt{temperature\_celsius} & celsius \\
      weather & \texttt{humidity\_percent} & percent \\
      air\_quality & \texttt{co2\_ppm} & ppm \\
      air\_quality & \texttt{particulate\_matter\_ug\_m3} & ug/m3 \\
      air\_quality & \texttt{air\_quality\_index} & aqi \\
      lamppost & \texttt{luminosity\_percent} & percent \\
      lamppost & \texttt{energy\_consumption\_kwh} & kWh \\
      lamppost & \texttt{light\_on} & boolean \\
      \bottomrule
    \end{tabular}
  \end{table}
  \begin{block}{Observacao}
    \texttt{light\_on} e armazenado como valor numerico booleano: \texttt{1.0} para ligado e \texttt{0.0} para desligado.
  \end{block}
\end{frame}

\begin{frame}{Historico Persistente no Gateway}
  \begin{itemize}
    \item O Gateway armazena as leituras em memoria para consultas durante a execucao.
    \item As leituras tambem sao persistidas em CSV para preservar historico entre reinicializacoes.
  \end{itemize}
\end{frame}

\begin{frame}{Parte 3: Cliente Web Analitico}
  \begin{itemize}
    \item Cliente implementado como aplicacao web em Python com Flask.
    \item O navegador acessa a interface HTTP na porta \texttt{5000}.
    \item O backend Flask usa \texttt{gateway\_api.py} para falar com o Gateway por TCP.
    \item A comunicacao Flask-Gateway usa \texttt{ClientRequest} e \texttt{ClientResponse} em Protobuf.
    \item A tela atualiza os dados periodicamente e apresenta graficos de serie temporal.
  \end{itemize}

  \begin{block}{Framing TCP}
    Cada mensagem TCP usa 4 bytes iniciais com o tamanho do payload, seguidos pelos bytes Protobuf.
  \end{block}
\end{frame}

\begin{frame}{Funcionalidades do Cliente Web}
  \begin{itemize}
    \item Listagem de fontes descobertas e respectivos estados.
    \item Listagem das leituras recentes recebidas pelo Gateway.
    \item Filtro por fonte e metrica para visualizacao no grafico.
    \item Consultas agregadas com \texttt{AVG}, \texttt{STDDEV}, \texttt{MIN} e \texttt{MAX}.
    \item Painel de controle do \texttt{lamppost}: ligar, desligar, consultar, ajustar luminosidade e simular falha.
    \item Atualizacao automatica da interface a cada 5 segundos.
  \end{itemize}
\end{frame}

\begin{frame}[fragile]{Mensagens Cliente-Gateway}
\footnotesize
  A interface web conversa por HTTP com o backend Flask. O backend traduz as operacoes da tela para mensagens TCP/Protobuf enviadas ao Gateway.
\vspace{0.15cm}
\begin{lstlisting}[style=proto,basicstyle=\ttfamily\tiny]
message ClientRequest {
  oneof request {
    ListSourcesRequest list_sources = 1;
    ListReadingsRequest list_readings = 2;
    SendCommandRequest send_command = 3;
    AggregateRequest aggregate = 4;
  }
}

message ClientResponse {
  bool success = 1;
  string message = 2;

  oneof response {
    ListSourcesResponse list_sources = 3;
    ListReadingsResponse list_readings = 4;
    SendCommandResponse send_command = 5;
    AggregateResponse aggregate = 6;
  }
}
\end{lstlisting}
\end{frame}

\begin{frame}[fragile]{Listagem de Fontes e Leituras}
\begin{lstlisting}[style=proto]
message SourceInfo {
  string name = 1;
  string address = 2;
  string ip = 3;
  int32 port = 4;
  bool controllable = 5;
  string status = 6;
  int64 last_seen_unix_ms = 7;
  string source_type = 8;
}

message ReadingInfo {
  string source_name = 1;
  string source_type = 2;
  string metric = 3;
  double value = 4;
  string unit = 5;
  int64 timestamp_unix_ms = 6;
}
\end{lstlisting}
\end{frame}

\begin{frame}{Controle do Lamppost}
  \begin{itemize}
    \item O \texttt{lamppost} é a unica fonte controlavel.
    \item O Gateway envia comandos por TCP para a porta \texttt{10002}.
    \item A fonte responde com status operacional, luminosidade, consumo e estado da luz.
    \item A fonte tem como atributos configuraveis a luminosidade e o estado da luz.
    \item O consumo de energia acumulado e exibido a partir das leituras mais recentes.
  \end{itemize}

\end{frame}

\begin{frame}[fragile]{Mensagens de Controle do Lamppost}
\begin{lstlisting}[style=proto]
message LamppostCommand {
  oneof command {
    LamppostTurnOnRequest turn_on = 1;
    LamppostTurnOffRequest turn_off = 2;
    LamppostGetStateRequest get_state = 3;
    LamppostSimulateFailureRequest simulate_failure = 4;
    LamppostSetLuminosityRequest set_luminosity = 5;
  }
}

message LamppostResponse {
  bool success = 1;
  string message = 2;
  bool active = 3;
  string status = 4;
  double luminosity_percent = 5;
  double energy_consumption_kwh = 6;
  bool light_on = 7;
}
\end{lstlisting}
\end{frame}

\begin{frame}[fragile]{Retorno de Comandos ao Cliente}
\begin{lstlisting}[style=proto]
message SendCommandResponse {
  bool success = 1;
  string message = 2;
  string source_status = 3;
  double luminosity_percent = 4;
  bool light_on = 5;
}
\end{lstlisting}

  \begin{itemize}
    \item O cliente exibe luminosidade e estado da luz apenas quando o comando tem sucesso.
    \item Em caso de falha local ou erro TCP, a resposta mostra apenas resultado, mensagem e status.
  \end{itemize}
\end{frame}

\begin{frame}{Consultas Agregadas}
  \begin{itemize}
    \item As consultas sao feitas por fonte, metrica, operacao e janela temporal.
    \item Operacoes suportadas: \texttt{AVG (Média)}, \texttt{STDDEV (Desvio Padrão)}, \texttt{MIN (Mínimo)} e \texttt{MAX (Máximo)}.
    \item O Gateway filtra leituras por \texttt{source\_name (Nome da fonte)}, \texttt{metric (Métrica)} e \texttt{window\_seconds (Intervalo de tempo em segundos)}.
    \item A resposta inclui valor, unidade, quantidade de amostras e janela usada.
    \item Quando \texttt{window\_seconds = 0}, todo o historico disponivel e considerado.
  \end{itemize}
\end{frame}

\begin{frame}[fragile]{Formato das Consultas}
\begin{lstlisting}[style=proto,basicstyle=\ttfamily\tiny]
enum AggregateOperation {
  AGGREGATE_OPERATION_AVG = 0;
  AGGREGATE_OPERATION_STDDEV = 1;
  AGGREGATE_OPERATION_MIN = 2;
  AGGREGATE_OPERATION_MAX = 3;
}

message AggregateRequest {
  string source_name = 1;
  string metric = 2;
  AggregateOperation operation = 3;
  int64 window_seconds = 4;
}

message AggregateResponse {
  string source_name = 1;
  string metric = 2;
  AggregateOperation operation = 3;
  double value = 4;
  int32 sample_count = 5;
  int64 window_seconds = 6;
  string unit = 7;
}
\end{lstlisting}
\end{frame}

\begin{frame}{Simulacao de Falhas}
  \small
  \begin{itemize}
    \item A mensagem \texttt{simulate\_failure} é enviado ao \texttt{lamppost}.
    \item A fonte responde ao comando e entra em falha temporaria por 15 segundos.
    \item Durante a falha, o \texttt{lamppost} nao responde discovery, nao envia leituras e nao responde mensagens TCP.
    \item O Gateway marca como \texttt{OFFLINE} quando a fonte nao responde a uma rodada de descoberta.
    \item Se um comando TCP for enviado durante a falha, o Gateway tambem trata o erro como indisponibilidade da fonte.
  \end{itemize}
\end{frame}

\begin{frame}{Tecnologias Utilizadas}
  \begin{table}
    \centering
    \begin{tabular}{ll}
      \toprule
      Tecnologia & Uso \\
      \midrule
      Go & Gateway, estado central, UDP/TCP, analytics e CSV \\
      Python & Fontes de dados e backend do cliente \\
      Flask & Servidor HTTP do cliente web \\
      HTML/CSS/JavaScript & Interface, atualizacao automatica e grafico \\
      Protocol Buffers & Contratos e serializacao das mensagens \\
      UDP multicast & Descoberta de fontes \\
      UDP unicast & Envio de leituras ao Gateway \\
      TCP & Flask-Gateway e controle do lamppost \\
      CSV & Persistencia do historico de leituras \\
      \bottomrule
    \end{tabular}
  \end{table}
\end{frame}

\begin{frame}{Organizacao do Codigo}
  \begin{itemize}
    \item \texttt{contracts/}: arquivos \texttt{.proto} compartilhados.
    \item \texttt{gateway/}: descoberta, estado, leituras, cliente TCP, analytics, historico e controle.
    \item \texttt{weather-sensor/}: fonte de temperatura e umidade.
    \item \texttt{air-quality-sensor/}: fonte de qualidade do ar.
    \item \texttt{lamppost/}: Poste de luz, fonte controlavel.
    \item \texttt{client/web\_client.py}: servidor Flask do cliente web.
    \item \texttt{client/gateway\_api.py}: camada TCP/Protobuf entre Flask e Gateway.
    \item \texttt{client/templates/} e \texttt{client/static/}: interface web.
  \end{itemize}
\end{frame}

\begin{frame}{Conclusao}
  \begin{itemize}
    \item O sistema integra processos independentes com comunicacao UDP e TCP.
    \item Protocol Buffers padroniza as mensagens entre Go e Python.
    \item O Gateway centraliza estado, leituras, historico, controle e consultas agregadas.
    \item O cliente web permite visualizar series temporais, consultar agregados e controlar a fonte atuadora.
    \item A fonte controlavel permite demonstrar comandos, estado operacional e falhas temporarias.
  \end{itemize}
\end{frame}

\end{document}
